\documentclass[12pt]{article}

\usepackage[a4paper, total={6in, 8in}]{geometry}
\usepackage{textcomp}

\begin{document}
\setlength{\parindent}{0pt}

\def \t {\textrightarrow}
\def \v {\vspace{1em}}
\def \bi {\begin{itemize}}
\def \ei {\end{itemize}}
\def \s[#1] {\section*{#1}}
\def \ss[#1] {\subsection*{#1}}
\def \sss[#1] {\subsubsection*{#1}}

\s[Montale]

\ss[Ti libero la fronte dai ghiaccioli]
Da "Le occasioni" \\
Clizia, dalle metamorfosi, significa girasole \t lei pero significa di piu per Montale di quanto Eleonora Duse abbia significato per D annunzio \\
Montale afferma che ci sono a volte delle occasioni \t sono opportunita che arrivano inaspettatamente \t arrivano come una illuminazione, e Clizia è la possiblita dell alba, illuminazione, redenzione \\
Clizia era dantista \t ma clizia gli viene allontana con le leggi razziali \\
Si coglie l'aspetto di una enigmaticita che caratterizza lo stile de "Le occasioni" \\
Le visiting angels, arrivano a salvare il poeta \t hanno davvero potere consolatorio?

\v

Due strofe, un tono connoquiale (ti libero la fronte) \t ghiaccioli è il freddo \t la figura arriva da dove? non si sa \t ghiaccioli porano a riflettere sull aridita e l impossibilita \\
Alte nebulose \t enj. \t penne lacerate dai cicloni

\v

Mezzodi \t immagini lapidarie, in un processo di agglomerazione \t sempre il "tu" con tono colloquiale, e presente il passato remoto che segna distanza \\
Torna l'immgine consueta \t il quadrante della sabbia del tempo \\
Giardino, contesto di limite \t il nespolo crea ombra \t il nespolo è anche albero della primavera, rappresenta la bella stagione \\
Ombra nera \\
Iperbato / anastrofe \t e enj. che separa sole e freddoloso \t anche le altre ombre si traggono, inconscie pero \t questo esprime il rapporto privilegiato con la visiting angels (ovvero clizia) \\
La comunicazoine con Clizia è solo loro \t le altre ombre non sanno che è qui \\
Con una presenza che non si cancella anche nell assenza \\
La secchezza di questo stile \t non c'è nulla di d annunziano, tranne il nespolo

\ss[Lo sai debbo riperdeti e non posso]
Il testo puo essere considerato molto toccante \t il contesto sotteso è l allontamento di clizia \\
Pirandello con Marta Abba, mentre Carducci con Piva \\
Si coglie il dolore, il distacco e l abbandono \\
"Cerco il segno smarrito" \t il pegno è il dono, che ebbi da te \\
Versi sono enigmatici \t ci sono termini che riprendono la desolazione interiore, come lo spiro salino, primavera che diventa scura \\
Polvere della sera \t ronzio con onomatopea \t lo spazio che evoca l immagine di un unghia che strazia i vetri \t a meta del verso 10 esplicita il suo bisogno profondo \\
Il pegno è un patto tacito o esplicito tra due persone \t il non trovare l anello che non tiene, l arco, porta alla paradigamtica espressione dell ultimo verso \\
Brandeis era una dantista \t ripresa qua \t Contini dice "Dante poeta concentrico"
\end{document}
